\documentclass[12pt,a4paper]{article}

\usepackage[T2A]{fontenc}
\usepackage[utf8]{inputenc}
\usepackage[russian]{babel}
\usepackage{amsmath,amssymb,amsthm}
\usepackage{mathtools}
\usepackage{algorithm}
\usepackage{algpseudocode}
\usepackage{array}
\usepackage{booktabs}
\usepackage{tikz}
\usepackage{enumitem}
\usepackage{hyperref}
\usepackage{geometry}
\geometry{margin=2.2cm}

\usetikzlibrary{arrows.meta,positioning,calc}

\newtheorem{definition}{Определение}
\newtheorem{theorem}{Теорема}
\newtheorem{lemma}{Лемма}
\newtheorem{remark}{Замечание}
\newtheorem{example}{Пример}

\DeclareMathOperator{\defop}{def}
\DeclareMathOperator{\dom}{dom}
\DeclareMathOperator{\useop}{use}
\DeclareMathOperator{\IN}{IN}
\DeclareMathOperator{\OUT}{OUT}
\DeclareMathOperator{\GEN}{GEN}
\DeclareMathOperator{\KILL}{KILL}
\DeclareMathOperator{\pred}{pred}
\DeclareMathOperator{\succs}{succ}
\DeclareMathOperator{\Dom}{Dom}
\DeclareMathOperator{\idom}{idom}
\DeclareMathOperator{\DF}{DF}
\DeclareMathOperator{\LiveIn}{LiveIn}
\DeclareMathOperator{\LiveOut}{LiveOut}

\title{Билет №38 \\ \small Оптимизация программ при их компиляции. Оптимизация базовых блоков, чистка циклов, распределение регистров. Анализ графов потока управления и потока данных. (ИТМО) \\ Оптимизация программ при их компиляции. Оптимизация базовых блоков, чистка циклов. Анализ графов потока управления и потока данных. Отношение доминирования и его свойства, построение границы области доминирования вершины, выделение сильно связанных компонент графа. Построение графа зависимостей. Перевод программы в SSA- представление и обратно. Глобальная и межпроцедурная оптимизация. (СпБГУ)}
\author{Сериков Владислав Александрович}
\date{13 мая 2026}

\begin{document}

\maketitle
\tableofcontents
\newpage

\section{Общая схема оптимизирующего компилятора}

Оптимизирующий компилятор обычно включает следующие фазы:
\begin{enumerate}
    \item лексический, синтаксический и семантический анализ;
    \item построение промежуточного представления программы;
    \item анализ потока управления;
    \item анализ потока данных;
    \item локальные, глобальные, цикловые и межпроцедурные оптимизации;
    \item распределение регистров;
    \item генерация целевого кода;
    \item машинно-зависимые оптимизации.
\end{enumerate}

\begin{definition}
\textbf{Оптимизация программы} --- преобразование программы, сохраняющее её наблюдаемую семантику, но улучшающее некоторые характеристики: время исполнения, объём кода, потребление памяти, энергопотребление, число обращений к памяти и т.д.
\end{definition}

Оптимизации делятся на:
\begin{itemize}
    \item \textbf{локальные} --- внутри одного базового блока;
    \item \textbf{глобальные} --- внутри одной процедуры;
    \item \textbf{цикловые} --- использующие структуру циклов;
    \item \textbf{межпроцедурные} --- анализирующие взаимодействие нескольких процедур;
    \item \textbf{машинно-независимые} --- выполняемые над абстрактным промежуточным представлением;
    \item \textbf{машинно-зависимые} --- учитывающие архитектуру процессора.
\end{itemize}

\section{Промежуточное представление программы}

Оптимизации обычно выполняются не над исходным текстом и не сразу над машинным кодом, а над промежуточным представлением.

Типичные формы промежуточного представления:
\begin{itemize}
    \item трёхадресный код;
    \item статическое однократное присваивание, SSA;
    \item деревья выражений;
    \item графы зависимостей;
    \item низкоуровневые IR, близкие к машинному коду.
\end{itemize}

\subsection{Трёхадресный код}

\begin{definition}
\textbf{Трёхадресный код} --- линейное промежуточное представление, в котором каждая команда содержит не более трёх адресов: два операнда и результат.
\end{definition}

Примеры команд:
\[
x := y + z,
\qquad
x := -y,
\qquad
x := y,
\qquad
\textbf{if } x < y \textbf{ goto } L,
\qquad
\textbf{goto } L.
\]

Например выражение
\[
a = b * c + d * e
\]
может быть представлено так:
\[
\begin{array}{l}
t_1 := b * c,\\
t_2 := d * e,\\
a := t_1 + t_2.
\end{array}
\]

\section{Базовые блоки и граф потока управления}

\subsection{Базовый блок}

\begin{definition}
\textbf{Базовый блок} --- максимальная последовательность команд с единственной точкой входа и единственной точкой выхода: управление входит только в первую команду блока и, если вошло, проходит последовательно до последней команды, не переходя внутрь других блоков.
\end{definition}

\subsection{Лидеры базовых блоков}

Для разбиения линейного кода на базовые блоки используется понятие лидеров.

\begin{definition}
\textbf{Лидер} --- первая команда некоторого базового блока.
\end{definition}

Лидерами являются:
\begin{enumerate}
    \item первая команда программы или процедуры;
    \item цель любого условного или безусловного перехода;
    \item команда, непосредственно следующая за условным или безусловным переходом.
\end{enumerate}

\subsection{Алгоритм построения базовых блоков}

\begin{algorithm}[H]
\caption{Построение базовых блоков}
\begin{algorithmic}[1]
\Require Последовательность команд трёхадресного кода
\Ensure Множество базовых блоков
\State Пометить первую команду как лидера.
\For{каждая команда перехода}
    \State Пометить цель перехода как лидера.
    \State Если после перехода есть следующая команда, пометить её как лидера.
\EndFor
\State Для каждого лидера построить блок из команд от него до команды перед следующим лидером.
\end{algorithmic}
\end{algorithm}

\begin{example}
Рассмотрим код:
\[
\begin{array}{rcl}
1&:& i := 0\\
2&:& s := 0\\
3&:& \textbf{if } i \geq n \textbf{ goto } 8\\
4&:& t := a[i]\\
5&:& s := s + t\\
6&:& i := i + 1\\
7&:& \textbf{goto } 3\\
8&:& \textbf{return } s
\end{array}
\]

Лидеры: 1, 3, 4, 8. Тогда базовые блоки:
\[
\begin{array}{ll}
B_1: & 1,2,\\
B_2: & 3,\\
B_3: & 4,5,6,7,\\
B_4: & 8.
\end{array}
\]
\end{example}

\subsection{Граф потока управления}

\begin{definition}
\textbf{Граф потока управления}, CFG, --- ориентированный граф
\[
G = (V,E),
\]
где вершины $V$ --- базовые блоки, а дуга $(B_i,B_j)\in E$ существует, если управление может непосредственно перейти из блока $B_i$ в блок $B_j$.
\end{definition}

Для предыдущего примера CFG имеет вид:

\[
\begin{tikzpicture}[
    node distance=1.7cm,
    block/.style={draw, rounded corners, minimum width=1.6cm, minimum height=0.8cm},
    >=Stealth
]
\node[block] (B1) {$B_1$};
\node[block, below of=B1] (B2) {$B_2$};
\node[block, below left=1.5cm and 1.6cm of B2] (B3) {$B_3$};
\node[block, below right=1.5cm and 1.6cm of B2] (B4) {$B_4$};

\draw[->] (B1) -- (B2);
\draw[->] (B2) -- node[left] {ложь} (B3);
\draw[->] (B2) -- node[right] {истина} (B4);
\draw[->] (B3) |- (B2);
\end{tikzpicture}
\]

\section{Оптимизация базовых блоков}

Оптимизация внутри базового блока называется локальной. Она не требует анализа всех путей программы и обычно выполняется быстрее глобальных оптимизаций.

\subsection{DAG базового блока}

Для локальной оптимизации часто строят ориентированный ациклический граф выражений.

\begin{definition}
\textbf{DAG базового блока} --- ориентированный ациклический граф, в котором листья соответствуют исходным значениям переменных и константам, а внутренние вершины --- операциям. С каждой вершиной связывается множество имён переменных, содержащих вычисленное в этой вершине значение.
\end{definition}

DAG позволяет:
\begin{itemize}
    \item устранять общие подвыражения;
    \item устранять лишние присваивания;
    \item выполнять свёртку констант;
    \item выявлять мёртвые вычисления внутри блока;
    \item упрощать алгебраические выражения.
\end{itemize}

\subsection{Алгоритм построения DAG базового блока}

\begin{algorithm}[H]
\caption{Построение DAG для базового блока}
\begin{algorithmic}[1]
\Require Последовательность команд базового блока
\Ensure DAG выражений
\For{каждая команда $x := y \ \mathrm{op}\ z$}
    \State Найти или создать вершины для текущих значений $y$ и $z$.
    \State Если уже существует вершина операции $\mathrm{op}$ с этими потомками, использовать её.
    \State Иначе создать новую вершину операции.
    \State Удалить имя $x$ из старой вершины, если оно было связано с ней.
    \State Добавить имя $x$ к найденной или созданной вершине.
\EndFor
\end{algorithmic}
\end{algorithm}

\begin{example}
Базовый блок:
\[
\begin{array}{l}
a := b + c,\\
d := b + c,\\
e := a + d.
\end{array}
\]

Команды $a := b+c$ и $d := b+c$ вычисляют одно и то же выражение. После построения DAG видно, что второе сложение можно заменить копированием:
\[
\begin{array}{l}
a := b + c,\\
d := a,\\
e := a + d.
\end{array}
\]
\end{example}

\subsection{Устранение общих подвыражений}

\begin{definition}
Выражение $x \ \mathrm{op}\ y$ является \textbf{общим подвыражением}, если оно уже вычислялось ранее и значения $x$ и $y$ с тех пор не изменялись.
\end{definition}

Пример:
\[
\begin{array}{l}
t_1 := a + b,\\
t_2 := c + d,\\
t_3 := a + b.
\end{array}
\]

Если $a$ и $b$ не изменялись между первой и третьей строками, то можно заменить:
\[
t_3 := t_1.
\]

\subsection{Свёртка и распространение констант}

\begin{definition}
\textbf{Свёртка констант} --- вычисление выражений с известными константными операндами во время компиляции.
\end{definition}

Пример:
\[
x := 2 * 3 + 4
\]
заменяется на
\[
x := 10.
\]

\begin{definition}
\textbf{Распространение констант} --- замена использования переменной известным константным значением.
\end{definition}

Пример:
\[
\begin{array}{l}
x := 5,\\
y := x + 1
\end{array}
\quad \Longrightarrow \quad
\begin{array}{l}
x := 5,\\
y := 6.
\end{array}
\]

\subsection{Алгебраические упрощения}

Типичные преобразования:
\[
x + 0 \Rightarrow x,
\qquad
x - 0 \Rightarrow x,
\qquad
x * 1 \Rightarrow x,
\qquad
x / 1 \Rightarrow x,
\qquad
x * 0 \Rightarrow 0.
\]

Некоторые преобразования требуют осторожности. Например, для чисел с плавающей точкой преобразования могут быть недопустимы из-за переполнений, NaN, бесконечностей и неассоциативности арифметики:
\[
(a+b)+c \not\equiv a+(b+c)
\]
в арифметике с плавающей точкой.

\subsection{Удаление мёртвого кода}

\begin{definition}
Присваивание называется \textbf{мёртвым}, если вычисленное значение далее не используется ни на одном возможном пути исполнения до следующего присваивания той же переменной.
\end{definition}

Пример:
\[
\begin{array}{l}
x := a + b,\\
x := c + d,\\
y := x + 1.
\end{array}
\]

Первое присваивание переменной $x$ мёртвое, так как его результат перезаписывается до использования.

\subsection{Копировочные преобразования}

Команда
\[
x := y
\]
создаёт копию значения. Если далее используется $x$, можно попытаться заменить его на $y$:
\[
\begin{array}{l}
x := y,\\
z := x + 1
\end{array}
\quad \Longrightarrow \quad
\begin{array}{l}
x := y,\\
z := y + 1.
\end{array}
\]

Если после этого присваивание $x := y$ становится мёртвым, его можно удалить.

\section{Анализ потока данных}

\subsection{Общая постановка}

\begin{definition}
\textbf{Анализ потока данных} --- метод статического анализа программы, при котором для каждой точки программы вычисляется некоторая информация, справедливая для всех или некоторых путей исполнения.
\end{definition}

Классическая задача анализа потока данных задаётся:
\begin{itemize}
    \item графом потока управления $G=(V,E)$;
    \item направлением анализа: прямой или обратный;
    \item множеством фактов;
    \item оператором объединения фактов;
    \item передаточными функциями блоков.
\end{itemize}

Для каждого блока $B$ обычно вычисляются множества:
\[
\IN[B], \qquad \OUT[B].
\]

Для прямого анализа:
\[
\IN[B] = \bigwedge_{P \in \pred(B)} \OUT[P],
\qquad
\OUT[B] = f_B(\IN[B]).
\]

Для обратного анализа:
\[
\OUT[B] = \bigwedge_{S \in \succs(B)} \IN[S],
\qquad
\IN[B] = f_B(\OUT[B]).
\]

Оператор $\wedge$ может быть:
\begin{itemize}
    \item объединением множеств $\cup$;
    \item пересечением множеств $\cap$;
    \item операцией взятия наибольшей нижней грани в некоторой решётке.
\end{itemize}

\subsection{Решётки и монотонные функции}

\begin{definition}
Частично упорядоченное множество $(L,\sqsubseteq)$ называется \textbf{решёткой}, если для любых $a,b\in L$ существуют:
\begin{itemize}
    \item точная верхняя грань $a \sqcup b$;
    \item точная нижняя грань $a \sqcap b$.
\end{itemize}
\end{definition}

\begin{definition}
Функция $f:L\to L$ называется \textbf{монотонной}, если
\[
x \sqsubseteq y \Rightarrow f(x) \sqsubseteq f(y).
\]
\end{definition}

\begin{theorem}[Существование неподвижной точки для монотонного анализа на конечной решётке]
Пусть $L$ --- конечная решётка, а $F:L^n\to L^n$ --- монотонная функция. Тогда итерационный алгоритм, начинающийся с верхнего или нижнего элемента решётки и последовательно применяющий $F$, достигает неподвижной точки за конечное число шагов.
\end{theorem}

\begin{proof}
Так как $L$ конечна, то и декартово произведение $L^n$ конечно. Итерационный алгоритм строит последовательность
\[
x_0, x_1, x_2, \dots,
\qquad
x_{k+1}=F(x_k).
\]
В классических анализах инициализация и направление итераций выбираются так, что последовательность монотонна: либо
\[
x_0 \sqsubseteq x_1 \sqsubseteq x_2 \sqsubseteq \dots,
\]
либо
\[
x_0 \sqsupseteq x_1 \sqsupseteq x_2 \sqsupseteq \dots.
\]
Поскольку множество $L^n$ конечно, в нём не может быть бесконечной строго возрастающей или строго убывающей цепи. Следовательно, существует номер $m$, начиная с которого
\[
x_m=x_{m+1}.
\]
Но $x_{m+1}=F(x_m)$, значит,
\[
F(x_m)=x_m.
\]
То есть $x_m$ является неподвижной точкой. Теорема доказана.
\end{proof}

\section{Классические задачи анализа потока данных}

\subsection{Достигающие определения}

\begin{definition}
Определение переменной $x$ в точке $d$ \textbf{достигает} точки $p$, если существует путь от $d$ до $p$, на котором $x$ не переопределяется.
\end{definition}

Это прямой анализ с объединением по предшественникам.

Для блока $B$:
\begin{itemize}
    \item $\GEN[B]$ --- определения, порождённые в $B$ и достигающие конца блока;
    \item $\KILL[B]$ --- определения тех же переменных, уничтожаемые определениями в $B$.
\end{itemize}

Уравнения:
\[
\IN[B] = \bigcup_{P\in \pred(B)} \OUT[P],
\]
\[
\OUT[B] = \GEN[B] \cup (\IN[B] \setminus \KILL[B]).
\]

\begin{example}
Пусть:
\[
\begin{array}{l}
d_1: x := 1,\\
d_2: y := x + 1,\\
d_3: x := 2,\\
d_4: z := x + y.
\end{array}
\]
Определение $d_1$ достигает $d_2$, но не достигает $d_4$, потому что между ними есть переопределение $x$ в $d_3$. Определение $d_3$ достигает $d_4$.
\end{example}

\subsection{Живые переменные}

\begin{definition}
Переменная $x$ является \textbf{живой} в точке $p$, если существует путь из $p$ к некоторому использованию $x$, на котором $x$ до этого использования не переопределяется.
\end{definition}

Это обратный анализ с объединением по преемникам.

Для блока $B$:
\begin{itemize}
    \item $\useop[B]$ --- переменные, используемые в $B$ до первого определения в $B$;
    \item $\defop[B]$ --- переменные, определяемые в $B$.
\end{itemize}

Уравнения:
\[
\OUT[B] = \bigcup_{S\in \succs(B)} \IN[S],
\]
\[
\IN[B] = \useop[B] \cup (\OUT[B]\setminus \defop[B]).
\]

Живость используется:
\begin{itemize}
    \item для удаления мёртвого кода;
    \item для распределения регистров;
    \item для построения графа интерференции.
\end{itemize}

\subsection{Доступные выражения}

\begin{definition}
Выражение $e$ \textbf{доступно} в точке $p$, если на каждом пути от входа программы до $p$ выражение $e$ уже вычислялось, и после последнего вычисления ни один из его операндов не изменялся.
\end{definition}

Это прямой анализ с пересечением по предшественникам.

Уравнения:
\[
\IN[B] = \bigcap_{P\in\pred(B)} \OUT[P],
\]
\[
\OUT[B] = \GEN[B] \cup (\IN[B]\setminus \KILL[B]).
\]

Используется для глобального устранения общих подвыражений.

\subsection{Очень занятые выражения}

\begin{definition}
Выражение $e$ является \textbf{очень занятым} в точке $p$, если на каждом пути из $p$ оно будет вычислено до изменения любого из его операндов.
\end{definition}

Это обратный анализ с пересечением:
\[
\OUT[B] = \bigcap_{S\in\succs(B)} \IN[S],
\]
\[
\IN[B] = \GEN[B] \cup (\OUT[B]\setminus \KILL[B]).
\]

Используется при выносе вычислений вперёд.

\section{Отношение доминирования}

\subsection{Определение}

Пусть $G=(V,E)$ --- CFG с выделенной входной вершиной $s$.

\begin{definition}
Вершина $d$ \textbf{доминирует} вершину $n$, если каждый путь из входной вершины $s$ в $n$ проходит через $d$.
Обозначение:
\[
d \dom n.
\]
\end{definition}

Очевидно:
\[
s \dom n
\]
для любой достижимой вершины $n$.

\begin{definition}
Вершина $d$ \textbf{строго доминирует} $n$, если
\[
d \dom n
\quad \text{и} \quad
d\neq n.
\]
\end{definition}

\begin{definition}
Вершина $d$ называется \textbf{непосредственным доминатором} вершины $n$, если $d$ строго доминирует $n$ и не существует вершины $c$, такой что
\[
d \dom c,\qquad c \dom n,\qquad c\neq d,\qquad c\neq n.
\]
Обозначение:
\[
d = \idom(n).
\]
\end{definition}

\subsection{Свойства доминирования}

\begin{theorem}[Основные свойства доминирования]
Для достижимых вершин CFG отношение доминирования обладает следующими свойствами:
\begin{enumerate}
    \item рефлексивность: $n\dom n$;
    \item антисимметричность: если $a\dom b$ и $b\dom a$, то $a=b$;
    \item транзитивность: если $a\dom b$ и $b\dom c$, то $a\dom c$.
\end{enumerate}
Следовательно, доминирование является частичным порядком на достижимых вершинах.
\end{theorem}

\begin{proof}
Докажем каждое свойство.

\textbf{1. Рефлексивность.}
Любой путь из входа $s$ в вершину $n$ заканчивается в $n$, следовательно, проходит через $n$. Поэтому $n\dom n$.

\textbf{2. Антисимметричность.}
Пусть $a\dom b$ и $b\dom a$.
Предположим, что $a\neq b$.
Так как $a\dom b$, каждый путь из $s$ в $b$ проходит через $a$.
Так как $b\dom a$, каждый путь из $s$ в $a$ проходит через $b$.

Рассмотрим любой путь из $s$ в $a$. Он должен пройти через $b$ до достижения $a$.
Но путь из $s$ в $b$ должен пройти через $a$ до достижения $b$.
Следовательно, на конечном простом пути до первого достижения одной из этих вершин возникает необходимость пройти через другую раньше, что невозможно без цикла, содержащего уже достигнутую вершину.
Если удалить циклы из пути, получим простой путь, на котором это противоречие сохраняется.
Значит, предположение $a\neq b$ неверно, и $a=b$.

\textbf{3. Транзитивность.}
Пусть $a\dom b$ и $b\dom c$.
Возьмём произвольный путь из $s$ в $c$.
Так как $b\dom c$, этот путь проходит через $b$.
Рассмотрим его префикс от $s$ до первого вхождения $b$.
Так как $a\dom b$, этот префикс проходит через $a$.
Значит, исходный путь из $s$ в $c$ также проходит через $a$.
Поскольку путь был произвольным, $a\dom c$.

Все три свойства доказаны.
\end{proof}

\subsection{Множества доминаторов}

Для каждой вершины $n$ обозначим множество её доминаторов:
\[
\Dom(n)=\{d\in V \mid d\dom n\}.
\]

Для входной вершины:
\[
\Dom(s)=\{s\}.
\]

Для остальных вершин:
\[
\Dom(n)=\{n\}\cup \bigcap_{p\in \pred(n)} \Dom(p).
\]

\begin{theorem}[Корректность уравнений доминирования]
Для достижимой вершины $n\neq s$ справедливо:
\[
\Dom(n)=\{n\}\cup \bigcap_{p\in \pred(n)} \Dom(p).
\]
\end{theorem}

\begin{proof}
Докажем включения в обе стороны.

\textbf{1. Пусть $d\in \Dom(n)$.}
Тогда каждый путь из $s$ в $n$ проходит через $d$.
Если $d=n$, то $d\in \{n\}$.
Пусть теперь $d\neq n$.
Любой путь из $s$ к предшественнику $p\in\pred(n)$ можно продолжить дугой $(p,n)$ до пути из $s$ в $n$.
Так как $d\neq n$ и $d$ доминирует $n$, вершина $d$ должна встретиться на участке пути до $p$.
Следовательно, $d$ доминирует каждого предшественника $p$.
Значит,
\[
d\in \bigcap_{p\in\pred(n)} \Dom(p).
\]

\textbf{2. Пусть}
\[
d\in \{n\}\cup \bigcap_{p\in\pred(n)} \Dom(p).
\]
Если $d=n$, то $d\dom n$ по рефлексивности.
Иначе $d$ доминирует каждого предшественника $p$ вершины $n$.
Любой путь из $s$ в $n$ перед последним шагом попадает в некоторого предшественника $p$ вершины $n$.
Поскольку $d\dom p$, этот путь до $p$ проходит через $d$.
Следовательно, весь путь до $n$ проходит через $d$.
Значит, $d\dom n$.

Итак, множества совпадают.
\end{proof}

\subsection{Итерационный алгоритм вычисления доминаторов}

\begin{algorithm}[H]
\caption{Итерационное вычисление множеств доминаторов}
\begin{algorithmic}[1]
\Require CFG $G=(V,E)$, входная вершина $s$
\Ensure $\Dom(n)$ для всех $n\in V$
\State $\Dom(s) := \{s\}$
\For{каждая вершина $n\neq s$}
    \State $\Dom(n) := V$
\EndFor
\Repeat
    \State $changed := false$
    \For{каждая вершина $n\neq s$}
        \State $new := \{n\} \cup \bigcap_{p\in\pred(n)} \Dom(p)$
        \If{$new \neq \Dom(n)$}
            \State $\Dom(n):=new$
            \State $changed := true$
        \EndIf
    \EndFor
\Until{$changed=false$}
\end{algorithmic}
\end{algorithm}

\begin{example}
Пусть CFG:

\[
\begin{tikzpicture}[
    node distance=1.5cm,
    block/.style={draw, circle, minimum size=0.8cm},
    >=Stealth
]
\node[block] (A) {$A$};
\node[block, below left of=A] (B) {$B$};
\node[block, below right of=A] (C) {$C$};
\node[block, below right of=B] (D) {$D$};
\node[block, below of=D] (E) {$E$};

\draw[->] (A) -- (B);
\draw[->] (A) -- (C);
\draw[->] (B) -- (D);
\draw[->] (C) -- (D);
\draw[->] (D) -- (E);
\end{tikzpicture}
\]

Тогда:
\[
\begin{array}{c|c}
n & \Dom(n)\\
\hline
A & \{A\}\\
B & \{A,B\}\\
C & \{A,C\}\\
D & \{A,D\}\\
E & \{A,D,E\}
\end{array}
\]
\end{example}

\subsection{Дерево доминаторов}

\begin{definition}
\textbf{Дерево доминаторов} --- дерево, в котором родителем каждой вершины $n\neq s$ является её непосредственный доминатор $\idom(n)$.
\end{definition}

Для предыдущего примера:
\[
\idom(B)=A,\qquad
\idom(C)=A,\qquad
\idom(D)=A,\qquad
\idom(E)=D.
\]

\section{Граница области доминирования}

\subsection{Определение}

\begin{definition}
\textbf{Граница доминирования} вершины $x$ --- множество вершин $y$, таких что $x$ доминирует некоторого непосредственного предшественника $y$, но не строго доминирует саму вершину $y$:
\[
\DF(x)=\{y\in V \mid \exists p\in\pred(y): x\dom p \ \text{и}\ x \not\dom_{\mathrm{strict}} y\}.
\]
\end{definition}

Интуитивно, $\DF(x)$ --- это множество точек, где пути, часть которых проходила через $x$, могут сливаться с путями, не находящимися полностью под доминированием $x$.

Граница доминирования используется при построении SSA-формы для определения мест вставки $\varphi$-функций.

\subsection{Пример границы доминирования}

Рассмотрим граф:

\[
\begin{tikzpicture}[
    node distance=1.4cm,
    block/.style={draw, circle, minimum size=0.8cm},
    >=Stealth
]
\node[block] (A) {$A$};
\node[block, below left=1.3cm and 1.5cm of A] (B) {$B$};
\node[block, below right=1.3cm and 1.5cm of A] (C) {$C$};
\node[block, below right=1.3cm and 1.5cm of B] (D) {$D$};

\draw[->] (A) -- (B);
\draw[->] (A) -- (C);
\draw[->] (B) -- (D);
\draw[->] (C) -- (D);
\end{tikzpicture}
\]

Вершина $A$ доминирует все вершины. Вершина $B$ доминирует саму себя, но не доминирует $D$, так как существует путь $A\to C\to D$.
При этом $B$ доминирует предшественника $D$, а именно $B$.
Следовательно:
\[
D\in \DF(B).
\]
Аналогично:
\[
D\in \DF(C).
\]

\subsection{Алгоритм построения границ доминирования}

Пусть уже построено дерево доминаторов.

\begin{algorithm}[H]
\caption{Построение границ доминирования}
\begin{algorithmic}[1]
\Require CFG $G=(V,E)$, непосредственные доминаторы $\idom$
\Ensure $\DF(n)$ для всех вершин $n$
\For{каждая вершина $n$}
    \State $\DF(n):=\varnothing$
\EndFor
\For{каждая вершина $y$, такая что $|\pred(y)|\geq 2$}
    \For{каждый $p\in \pred(y)$}
        \State $runner := p$
        \While{$runner \neq \idom(y)$}
            \State $\DF(runner):=\DF(runner)\cup\{y\}$
            \State $runner := \idom(runner)$
        \EndWhile
    \EndFor
\EndFor
\end{algorithmic}
\end{algorithm}

\begin{example}
Для ромбовидного графа:
\[
A\to B,\quad A\to C,\quad B\to D,\quad C\to D
\]
имеем:
\[
\idom(B)=A,\quad \idom(C)=A,\quad \idom(D)=A.
\]
Вершина $D$ имеет двух предшественников: $B$ и $C$.

Для предшественника $B$:
\[
runner=B,\qquad runner\neq \idom(D)=A,
\]
поэтому добавляем:
\[
D\in \DF(B).
\]

Для предшественника $C$ аналогично:
\[
D\in \DF(C).
\]

Итог:
\[
\DF(B)=\{D\},\qquad
\DF(C)=\{D\}.
\]
\end{example}

\section{Циклы в графе потока управления}

\subsection{Обратные дуги}

\begin{definition}
Дуга $(n,h)$ в CFG называется \textbf{обратной}, если вершина $h$ доминирует вершину $n$:
\[
h\dom n.
\]
\end{definition}

В этом случае $h$ называется заголовком цикла.

\subsection{Естественный цикл}

\begin{definition}
\textbf{Естественный цикл} обратной дуги $(n,h)$ --- множество вершин, состоящее из $h$, $n$ и всех вершин, из которых можно попасть в $n$, не проходя через $h$.
\end{definition}

\subsection{Алгоритм построения естественного цикла}

\begin{algorithm}[H]
\caption{Построение естественного цикла для обратной дуги $(n,h)$}
\begin{algorithmic}[1]
\Require CFG, обратная дуга $(n,h)$
\Ensure Множество вершин цикла $Loop$
\State $Loop := \{h,n\}$
\State $Stack := [n]$
\While{$Stack$ не пуст}
    \State $m := pop(Stack)$
    \For{каждый $p\in \pred(m)$}
        \If{$p\notin Loop$}
            \State добавить $p$ в $Loop$
            \State $push(Stack,p)$
        \EndIf
    \EndFor
\EndWhile
\end{algorithmic}
\end{algorithm}

\begin{example}
Пусть:
\[
A\to B,\quad B\to C,\quad C\to D,\quad D\to B.
\]
Если $B\dom D$, то $(D,B)$ --- обратная дуга. Естественный цикл:
\[
\{B,C,D\}.
\]
\end{example}

\section{Чистка циклов и оптимизации циклов}

Циклы являются важнейшей целью оптимизации, так как команды внутри циклов выполняются многократно.

\subsection{Инварианты цикла}

\begin{definition}
Выражение называется \textbf{инвариантным относительно цикла}, если его значение не меняется при различных итерациях цикла.
\end{definition}

Например:
\[
\begin{array}{l}
i := 0\\
\textbf{while } i<n \textbf{ do}\\
\quad t := a*b\\
\quad x[i] := t+i\\
\quad i := i+1
\end{array}
\]
Если $a$ и $b$ не изменяются внутри цикла, то $a*b$ является инвариантным выражением.

\subsection{Вынос инвариантов из цикла}

\begin{definition}
\textbf{Вынос инвариантов из цикла} --- преобразование, при котором вычисление инвариантного выражения переносится в предзаголовок цикла.
\end{definition}

Предзаголовок --- блок, который выполняется непосредственно перед входом в цикл и из которого управление переходит в заголовок цикла.

Исходный код:
\[
\begin{array}{l}
\textbf{while } i<n \textbf{ do}\\
\quad t := a*b\\
\quad x[i] := t+i\\
\quad i := i+1
\end{array}
\]

После оптимизации:
\[
\begin{array}{l}
t := a*b\\
\textbf{while } i<n \textbf{ do}\\
\quad x[i] := t+i\\
\quad i := i+1
\end{array}
\]

\subsection{Условия корректности выноса инварианта}

Команда
\[
x := y \ \mathrm{op}\ z
\]
может быть вынесена из цикла, если:
\begin{enumerate}
    \item операнды $y$ и $z$ либо определены вне цикла, либо сами являются результатами уже вынесенных инвариантных вычислений;
    \item команда доминирует все выходы из цикла, если её результат используется после цикла;
    \item переменная $x$ не переопределяется другим образом внутри цикла;
    \item вынос не меняет поведение при исключениях, переполнениях и побочных эффектах.
\end{enumerate}

\subsection{Алгоритм поиска инвариантных команд}

\begin{algorithm}[H]
\caption{Поиск инвариантных команд цикла}
\begin{algorithmic}[1]
\Require Цикл $L$
\Ensure Множество инвариантных команд
\State $Invariant := \varnothing$
\Repeat
    \State $changed := false$
    \For{каждая команда $s: x:=y \ \mathrm{op}\ z$ в $L$}
        \If{$s\notin Invariant$ и все определения $y,z$ находятся вне $L$ или в $Invariant$}
            \State добавить $s$ в $Invariant$
            \State $changed := true$
        \EndIf
    \EndFor
\Until{$changed=false$}
\end{algorithmic}
\end{algorithm}

\subsection{Индуктивные переменные}

\begin{definition}
\textbf{Индуктивная переменная} --- переменная, значение которой на каждой итерации цикла изменяется на постоянную величину или является линейной функцией другой индуктивной переменной.
\end{definition}

Пример:
\[
i := i+1
\]
задаёт базовую индуктивную переменную.

Если внутри цикла:
\[
j := 4*i+8,
\]
то $j$ является производной индуктивной переменной.

\subsection{Уменьшение силы операций}

\begin{definition}
\textbf{Уменьшение силы операций} --- замена дорогих операций более дешёвыми, например умножения сложением.
\end{definition}

Исходный код:
\[
\begin{array}{l}
\textbf{for } i=0 \textbf{ to } n-1:\\
\quad y := 4*i
\end{array}
\]

После оптимизации:
\[
\begin{array}{l}
y := 0\\
\textbf{for } i=0 \textbf{ to } n-1:\\
\quad \ldots\\
\quad y := y+4
\end{array}
\]

\subsection{Удаление лишних проверок}

Если компилятор доказывает, что индекс массива всегда находится в допустимом диапазоне, можно удалить проверку границ.

Пример:
\[
\begin{array}{l}
\textbf{for } i=0 \textbf{ to } n-1:\\
\quad a[i] := 0
\end{array}
\]
Если известно, что массив $a$ имеет длину $n$, проверка $0\leq i<n$ внутри цикла может быть вынесена или удалена.

\subsection{Разворачивание цикла}

\begin{definition}
\textbf{Разворачивание цикла} --- преобразование, уменьшающее число итераций за счёт дублирования тела цикла.
\end{definition}

Исходный код:
\[
\begin{array}{l}
\textbf{for } i=0 \textbf{ to } n-1:\\
\quad a[i] := a[i]+1
\end{array}
\]

Разворачивание с фактором 2:
\[
\begin{array}{l}
\textbf{for } i=0 \textbf{ to } n-2 \textbf{ step } 2:\\
\quad a[i] := a[i]+1\\
\quad a[i+1] := a[i+1]+1
\end{array}
\]

Плюсы:
\begin{itemize}
    \item меньше переходов;
    \item больше возможностей для планирования инструкций;
    \item возможна векторизация.
\end{itemize}

Минусы:
\begin{itemize}
    \item увеличение размера кода;
    \item ухудшение кэширования инструкций.
\end{itemize}

\section{Сильно связанные компоненты графа}

\subsection{Определение}

\begin{definition}
В ориентированном графе две вершины $u$ и $v$ называются \textbf{взаимно достижимыми}, если существует путь из $u$ в $v$ и путь из $v$ в $u$.
\end{definition}

\begin{definition}
\textbf{Сильно связная компонента}, SCC, --- максимальное по включению множество вершин ориентированного графа, в котором любые две вершины взаимно достижимы.
\end{definition}

SCC используются:
\begin{itemize}
    \item для анализа циклов в CFG;
    \item для оптимизации рекурсивных вызовов;
    \item для анализа зависимостей;
    \item при межпроцедурном анализе графа вызовов.
\end{itemize}

\subsection{Алгоритм Косарайю}

\begin{algorithm}[H]
\caption{Алгоритм Косарайю для поиска SCC}
\begin{algorithmic}[1]
\Require Ориентированный граф $G=(V,E)$
\Ensure Сильно связные компоненты
\State Выполнить DFS в $G$, записывая вершины в порядке завершения.
\State Построить транспонированный граф $G^T$, где все дуги обращены.
\State Рассматривать вершины в порядке убывания времени завершения DFS.
\State Для каждой ещё не посещённой вершины выполнить DFS в $G^T$.
\State Каждое дерево DFS во втором проходе является одной SCC.
\end{algorithmic}
\end{algorithm}

\begin{theorem}[Корректность алгоритма Косарайю]
Алгоритм Косарайю разбивает ориентированный граф на его сильно связные компоненты.
\end{theorem}

\begin{proof}
Сожмём каждую сильно связную компоненту исходного графа $G$ в одну вершину. Получим конденсацию графа $C(G)$. Этот граф ацикличен: если бы в нём существовал цикл между компонентами, то все вершины этих компонент были бы взаимно достижимы, а значит, они образовывали бы одну SCC, что противоречит максимальности компонент.

В первом DFS времена завершения обладают следующим свойством: если в конденсации есть дуга из компоненты $A$ в компоненту $B$, то максимальное время завершения вершины из $A$ больше максимального времени завершения вершины из $B$, если DFS сначала вошёл в $A$ и смог достичь $B$; если же DFS сначала посетил $B$, то из $B$ нельзя достичь $A$ в конденсации, поэтому $B$ завершится раньше, чем позднее начатый обход $A$. В любом случае вершины исходных компонент упорядочиваются так, что компоненты-источники в $C(G)$ имеют большие времена завершения.

Во втором проходе рассматривается транспонированный граф $G^T$. Источник в $C(G)$ становится стоком в $C(G^T)$. Если начать DFS из вершины компоненты с максимальным временем завершения, то в транспонированном графе из неё нельзя уйти в ещё не обработанную другую компоненту: иначе в исходной конденсации существовал бы путь, нарушающий порядок времён завершения. Следовательно, DFS во втором проходе обходит ровно одну SCC.

После удаления найденной компоненты рассуждение повторяется для оставшегося графа. Поэтому каждый запуск DFS во втором проходе выделяет ровно одну сильно связную компоненту, и все компоненты будут найдены.
\end{proof}

\subsection{Алгоритм Тарьяна}

Алгоритм Тарьяна находит SCC за один DFS.

Для каждой вершины $v$ хранятся:
\begin{itemize}
    \item $index[v]$ --- номер открытия вершины;
    \item $lowlink[v]$ --- минимальный номер открытия вершины, достижимой из $v$ по дереву DFS и не более чем одной обратной дуге;
    \item стек активных вершин.
\end{itemize}

\begin{algorithm}[H]
\caption{Алгоритм Тарьяна}
\begin{algorithmic}[1]
\State $index := 0$
\State $S :=$ пустой стек
\For{каждая вершина $v$}
    \If{$v$ ещё не посещена}
        \State $StrongConnect(v)$
    \EndIf
\EndFor

\Procedure{StrongConnect}{$v$}
    \State $index[v] := index$
    \State $lowlink[v] := index$
    \State $index := index+1$
    \State поместить $v$ в стек $S$
    \State $onStack[v] := true$
    \For{каждая дуга $(v,w)$}
        \If{$w$ ещё не посещена}
            \State $StrongConnect(w)$
            \State $lowlink[v] := \min(lowlink[v], lowlink[w])$
        \ElsIf{$onStack[w]$}
            \State $lowlink[v] := \min(lowlink[v], index[w])$
        \EndIf
    \EndFor
    \If{$lowlink[v]=index[v]$}
        \State начать новую SCC
        \Repeat
            \State $w := pop(S)$
            \State $onStack[w] := false$
            \State добавить $w$ в текущую SCC
        \Until{$w=v$}
    \EndIf
\EndProcedure
\end{algorithmic}
\end{algorithm}

Время работы:
\[
O(|V|+|E|).
\]

\section{Граф зависимостей}

\subsection{Зависимости между операциями}

\begin{definition}
\textbf{Граф зависимостей} --- ориентированный граф, вершины которого соответствуют операциям программы, а дуги задают ограничения на порядок выполнения операций.
\end{definition}

Основные виды зависимостей:

\begin{enumerate}
    \item \textbf{Истинная зависимость}, RAW, read after write:
    \[
    S_1: x:=\ldots,
    \qquad
    S_2: \ldots:=x+\ldots
    \]
    Операция $S_2$ должна выполняться после $S_1$.

    \item \textbf{Антизависимость}, WAR, write after read:
    \[
    S_1: \ldots:=x+\ldots,
    \qquad
    S_2: x:=\ldots
    \]
    Нельзя переместить $S_2$ перед $S_1$ без изменения значения, читаемого $S_1$.

    \item \textbf{Выходная зависимость}, WAW, write after write:
    \[
    S_1: x:=\ldots,
    \qquad
    S_2: x:=\ldots
    \]
    Порядок записей важен, если далее используется последнее значение.

    \item \textbf{Управляющая зависимость}:
    выполнение операции зависит от результата условного перехода.
\end{enumerate}

\subsection{Построение графа зависимостей}

\begin{algorithm}[H]
\caption{Построение простого графа зависимостей для последовательности команд}
\begin{algorithmic}[1]
\Require Последовательность команд
\Ensure Граф зависимостей
\State Создать вершину для каждой команды.
\For{каждая пара команд $S_i,S_j$, где $i<j$}
    \If{$S_i$ записывает переменную, которую $S_j$ читает}
        \State добавить дугу $S_i\to S_j$ типа RAW
    \EndIf
    \If{$S_i$ читает переменную, которую $S_j$ записывает}
        \State добавить дугу $S_i\to S_j$ типа WAR
    \EndIf
    \If{$S_i$ записывает переменную, которую $S_j$ записывает}
        \State добавить дугу $S_i\to S_j$ типа WAW
    \EndIf
\EndFor
\end{algorithmic}
\end{algorithm}

\begin{example}
Код:
\[
\begin{array}{l}
S_1: a := b+c,\\
S_2: d := a*2,\\
S_3: a := e+f.
\end{array}
\]

Зависимости:
\[
S_1 \to S_2 \quad \text{RAW по } a,
\]
\[
S_1 \to S_3 \quad \text{WAW по } a,
\]
\[
S_2 \to S_3 \quad \text{WAR по } a.
\]
\end{example}

\subsection{Граф зависимостей программы}

\begin{definition}
\textbf{Program Dependence Graph}, PDG, --- граф, объединяющий зависимости по данным и зависимости по управлению.
\end{definition}

PDG используется:
\begin{itemize}
    \item для слайсинга программ;
    \item для параллелизации;
    \item для обнаружения независимых вычислений;
    \item для оптимизаций с изменением порядка команд.
\end{itemize}

\section{SSA-представление}

\subsection{Определение SSA}

\begin{definition}
Программа находится в форме \textbf{статического однократного присваивания}, SSA, если каждая переменная определяется ровно один раз в тексте программы.
\end{definition}

Если переменная получает разные значения на разных путях управления, используется специальная $\varphi$-функция.

Пример:
\[
\begin{array}{l}
\textbf{if } c \textbf{ then}\\
\quad x := 1\\
\textbf{else}\\
\quad x := 2\\
y := x+3
\end{array}
\]

SSA-форма:
\[
\begin{array}{l}
\textbf{if } c \textbf{ then}\\
\quad x_1 := 1\\
\textbf{else}\\
\quad x_2 := 2\\
x_3 := \varphi(x_1,x_2)\\
y_1 := x_3+3.
\end{array}
\]

\subsection{Смысл $\varphi$-функции}

Команда
\[
x_3 := \varphi(x_1,x_2)
\]
означает: если управление пришло из первого предшественника блока, выбрать $x_1$; если из второго --- выбрать $x_2$.

$\varphi$-функция не является обычным вычислением во время выполнения; она описывает слияние значений на уровне промежуточного представления.

\section{Перевод программы в SSA}

Классический алгоритм состоит из двух больших этапов:
\begin{enumerate}
    \item вставка $\varphi$-функций;
    \item переименование переменных.
\end{enumerate}

\subsection{Вставка $\varphi$-функций}

Для каждой переменной $a$ рассматривается множество блоков, где она определяется:
\[
DefBlocks(a).
\]

$\varphi$-функции вставляются в итеративную границу доминирования этих блоков.

\begin{definition}
\textbf{Итеративная граница доминирования} множества $S$ определяется как наименьшее множество $IDF(S)$, содержащее $\DF(S)$ и замкнутое относительно повторного применения $\DF$:
\[
IDF(S)=\mu X.\left(\DF(S\cup X)\right).
\]
\end{definition}

\begin{algorithm}[H]
\caption{Вставка $\varphi$-функций}
\begin{algorithmic}[1]
\Require CFG, границы доминирования, множества $DefBlocks(a)$
\Ensure Программа с расставленными $\varphi$-функциями
\For{каждая переменная $a$}
    \State $WorkList := DefBlocks(a)$
    \State $HasAlready := \varnothing$
    \While{$WorkList$ не пуст}
        \State удалить некоторый блок $X$ из $WorkList$
        \For{каждый $Y\in \DF(X)$}
            \If{в $Y$ ещё нет $\varphi$ для $a$}
                \State вставить $a := \varphi(\ldots)$ в начало $Y$
                \State добавить $Y$ в $HasAlready$
                \If{$Y\notin DefBlocks(a)$}
                    \State добавить $Y$ в $WorkList$
                \EndIf
            \EndIf
        \EndFor
    \EndWhile
\EndFor
\end{algorithmic}
\end{algorithm}

\subsection{Переименование переменных}

Идея: при обходе дерева доминаторов поддерживается стек текущих версий каждой переменной.

\begin{algorithm}[H]
\caption{Переименование переменных в SSA}
\begin{algorithmic}[1]
\Require CFG с $\varphi$-функциями, дерево доминаторов
\Ensure SSA-форма
\State Для каждой переменной $a$: $stack[a]:=$ пустой стек, $counter[a]:=0$
\State Вызвать $Rename(entry)$

\Procedure{NewName}{$a$}
    \State $counter[a]:=counter[a]+1$
    \State $name:=a_{counter[a]}$
    \State $push(stack[a],name)$
    \State \Return $name$
\EndProcedure

\Procedure{Rename}{$B$}
    \State $Pushed :=$ пустой список
    \For{каждая $\varphi$-функция $a:=\varphi(\ldots)$ в $B$}
        \State заменить определяемую переменную $a$ на $NewName(a)$
        \State добавить $a$ в $Pushed$
    \EndFor

    \For{каждая обычная команда в $B$}
        \State заменить каждое использование $a$ на $top(stack[a])$
        \If{команда определяет $a$}
            \State заменить определение $a$ на $NewName(a)$
            \State добавить $a$ в $Pushed$
        \EndIf
    \EndFor

    \For{каждый преемник $S$ блока $B$}
        \State В $\varphi$-функциях блока $S$ заполнить аргумент, соответствующий ребру $B\to S$, значением $top(stack[a])$
    \EndFor

    \For{каждый сын $C$ блока $B$ в дереве доминаторов}
        \State $Rename(C)$
    \EndFor

    \For{каждая переменная $a$ из $Pushed$ в обратном порядке}
        \State $pop(stack[a])$
    \EndFor
\EndProcedure
\end{algorithmic}
\end{algorithm}

\subsection{Минимальный пример построения SSA}

Исходный код:
\[
\begin{array}{l}
B_1: x:=0\\
\quad \textbf{if } c \textbf{ goto } B_2 \textbf{ else } B_3\\[2mm]
B_2: x:=1\\
\quad \textbf{goto } B_4\\[2mm]
B_3: x:=2\\
\quad \textbf{goto } B_4\\[2mm]
B_4: y:=x+1
\end{array}
\]

В блоке $B_4$ сходятся два определения $x$, поэтому нужна $\varphi$-функция:
\[
\begin{array}{l}
B_1: x_1:=0\\
\quad \textbf{if } c \textbf{ goto } B_2 \textbf{ else } B_3\\[2mm]
B_2: x_2:=1\\
\quad \textbf{goto } B_4\\[2mm]
B_3: x_3:=2\\
\quad \textbf{goto } B_4\\[2mm]
B_4: x_4:=\varphi(x_2,x_3)\\
\quad y_1:=x_4+1
\end{array}
\]

Заметим, что $x_1$ здесь не используется, если оба пути после $B_1$ обязательно переопределяют $x$.

\section{Оптимизации в SSA}

SSA-форма упрощает многие оптимизации.

\subsection{Разреженное распространение констант}

В SSA у каждой переменной ровно одно определение, поэтому легче распространять значения по графу использований.

Пример:
\[
\begin{array}{l}
x_1:=2\\
y_1:=3\\
z_1:=x_1+y_1
\end{array}
\quad \Rightarrow \quad
z_1:=5.
\]

\subsection{Удаление мёртвого кода в SSA}

Если SSA-имя не используется и команда не имеет побочных эффектов, команду можно удалить.

Пример:
\[
\begin{array}{l}
x_1:=a+b\\
y_1:=c+d
\end{array}
\]
Если $x_1$ нигде не используется, команда $x_1:=a+b$ удаляется.

\subsection{Устранение копий}

Команда:
\[
x_1:=y_1
\]
может быть удалена заменой всех использований $x_1$ на $y_1$.

\section{Перевод из SSA обратно}

После оптимизаций SSA нужно преобразовать в обычный код без $\varphi$-функций.

\subsection{Наивное удаление $\varphi$-функций}

Команда:
\[
x := \varphi(a,b)
\]
в блоке $B$ заменяется вставкой копий в предшественники блока $B$:
\[
\begin{array}{l}
P_1:\ x:=a\\
P_2:\ x:=b
\end{array}
\]

Однако наивный подход может быть некорректен при параллельной семантике $\varphi$-функций.

\subsection{Параллельная семантика $\varphi$-функций}

Все $\varphi$-функции в начале блока считаются выполняющимися одновременно.

Пример:
\[
\begin{array}{l}
x_2:=\varphi(y_1,\ldots)\\
y_2:=\varphi(x_1,\ldots)
\end{array}
\]
Если вставить копии последовательно:
\[
x:=y;\quad y:=x,
\]
то второе присваивание использует уже изменённый $x$. Поэтому может потребоваться временная переменная.

Правильно:
\[
t:=x;\quad x:=y;\quad y:=t.
\]

\subsection{Алгоритм удаления $\varphi$-функций}

\begin{algorithm}[H]
\caption{Удаление $\varphi$-функций}
\begin{algorithmic}[1]
\Require SSA-программа
\Ensure Обычная программа без $\varphi$-функций
\For{каждый блок $B$}
    \For{каждая $\varphi$-функция $x:=\varphi(a_1,\ldots,a_k)$ в $B$}
        \For{каждый предшественник $P_i$ блока $B$}
            \State добавить копию $x:=a_i$ на ребро $P_i\to B$
        \EndFor
    \EndFor
\EndFor
\State Если ребро критическое, предварительно разбить его новым блоком.
\State Упорядочить параллельные копии с использованием временных переменных при циклах копирования.
\State Удалить все $\varphi$-функции.
\end{algorithmic}
\end{algorithm}

\begin{definition}
Ребро $P\to B$ называется \textbf{критическим}, если $P$ имеет более одного преемника, а $B$ имеет более одного предшественника.
\end{definition}

Критические рёбра разбивают, чтобы вставка копий на ребро не изменила поведение других переходов.

\section{Глобальная оптимизация}

\begin{definition}
\textbf{Глобальная оптимизация} --- оптимизация в пределах одной процедуры, использующая информацию о всём CFG этой процедуры.
\end{definition}

Основные глобальные оптимизации:
\begin{itemize}
    \item глобальное устранение общих подвыражений;
    \item глобальное распространение констант;
    \item удаление мёртвого кода;
    \item вынос инвариантов из циклов;
    \item частичное устранение избыточности;
    \item оптимизация переходов;
    \item распределение регистров.
\end{itemize}

\subsection{Глобальное устранение общих подвыражений}

Если выражение доступно в точке программы и операнды не изменялись, повторное вычисление заменяется использованием ранее вычисленного значения.

Пример:
\[
\begin{array}{l}
B_1: t_1:=a+b\\
\quad \textbf{if } c \textbf{ goto } B_2 \textbf{ else } B_3\\
B_2: x:=a+b\\
B_3: y:=a+b
\end{array}
\]

Если $a+b$ доступно в $B_2$ и $B_3$, можно использовать $t_1$.

\subsection{Частичное устранение избыточности}

\begin{definition}
Выражение является \textbf{частично избыточным}, если оно уже вычислено на некоторых, но не на всех путях к данной точке.
\end{definition}

Идея PRE: вставить вычисление на недостающих путях так, чтобы затем устранить повторное вычисление.

Пример:
\[
\begin{array}{l}
\textbf{if } c \textbf{ then}\\
\quad t:=a+b\\
x:=a+b
\end{array}
\]
Выражение $a+b$ частично избыточно перед присваиванием $x$. Можно вычислить его также в ветке, где оно не вычислялось, и заменить второе вычисление.

\section{Межпроцедурная оптимизация}

\begin{definition}
\textbf{Межпроцедурная оптимизация} --- оптимизация, использующая информацию о нескольких процедурах программы.
\end{definition}

Она строится на:
\begin{itemize}
    \item графе вызовов;
    \item анализе побочных эффектов;
    \item анализе алиасов;
    \item межпроцедурном распространении констант;
    \item встраивании функций;
    \item удалении неиспользуемых процедур;
    \item специализации процедур.
\end{itemize}

\subsection{Граф вызовов}

\begin{definition}
\textbf{Граф вызовов} --- ориентированный граф, вершины которого соответствуют процедурам, а дуга $f\to g$ означает, что процедура $f$ может вызвать процедуру $g$.
\end{definition}

Если граф вызовов содержит циклы, это соответствует рекурсии. SCC графа вызовов позволяют анализировать группы взаимно рекурсивных процедур.

\subsection{Встраивание функций}

\begin{definition}
\textbf{Встраивание функции}, inline expansion, --- замена вызова функции телом этой функции.
\end{definition}

Пример:
\[
\begin{array}{l}
f(x)=x+1,\\
y:=f(a)
\end{array}
\quad \Rightarrow \quad
y:=a+1.
\]

Преимущества:
\begin{itemize}
    \item устраняются накладные расходы вызова;
    \item появляются новые возможности для локальных и глобальных оптимизаций;
    \item возможно распространение констант через параметры.
\end{itemize}

Недостатки:
\begin{itemize}
    \item увеличивается размер кода;
    \item может ухудшиться кэширование инструкций.
\end{itemize}

\subsection{Межпроцедурное распространение констант}

Если известно, что функция всегда вызывается с некоторым константным параметром, можно специализировать её.

Исходно:
\[
\begin{array}{l}
f(x,y)=x+y,\\
z:=f(10,a).
\end{array}
\]

После специализации:
\[
f_{10}(y)=10+y,
\qquad
z:=f_{10}(a).
\]

\subsection{Анализ побочных эффектов}

Для функции $f$ можно определить:
\[
Read(f) = \text{множество глобальных объектов, которые } f \text{ читает},
\]
\[
Write(f) = \text{множество глобальных объектов, которые } f \text{ изменяет}.
\]

Если:
\[
Write(f)=\varnothing
\]
и результат функции зависит только от аргументов, то вызов можно рассматривать как чистый. Это позволяет:
\begin{itemize}
    \item устранять повторные вызовы;
    \item переставлять вызовы;
    \item удалять неиспользуемые вызовы.
\end{itemize}

\subsection{Анализ алиасов}

\begin{definition}
Два выражения-указателя являются \textbf{алиасами}, если они могут указывать на одну и ту же область памяти.
\end{definition}

Пример:
\[
\begin{array}{l}
*p := 1;\\
x := *q;
\end{array}
\]
Если $p$ и $q$ могут быть алиасами, нельзя свободно переставить эти операции.

Анализ алиасов бывает:
\begin{itemize}
    \item \textbf{may-alias}: могут ли два указателя ссылаться на один объект;
    \item \textbf{must-alias}: обязаны ли два указателя ссылаться на один объект.
\end{itemize}

\section{Распределение регистров}

\subsection{Постановка задачи}

\begin{definition}
\textbf{Распределение регистров} --- сопоставление переменных и временных значений программы физическим регистрам процессора.
\end{definition}

Так как число физических регистров ограничено, некоторые значения приходится хранить в памяти. Это называется spill.

\begin{definition}
\textbf{Spill} --- размещение значения во временной ячейке памяти вместо регистра.
\end{definition}

\subsection{Интервалы живости}

\begin{definition}
\textbf{Интервал живости} переменной --- участок программы от её определения до последнего использования.
\end{definition}

Если интервалы живости двух переменных пересекаются, они не могут находиться в одном регистре.

\subsection{Граф интерференции}

\begin{definition}
\textbf{Граф интерференции} --- неориентированный граф, вершины которого соответствуют переменным или временным значениям, а ребро между двумя вершинами означает, что соответствующие значения одновременно живы и не могут быть помещены в один регистр.
\end{definition}

Распределение регистров сводится к раскраске графа.

\begin{definition}
\textbf{$k$-раскраска графа} --- назначение каждой вершине одного из $k$ цветов так, что смежные вершины имеют разные цвета.
\end{definition}

Если доступно $k$ регистров, нужно раскрасить граф интерференции в $k$ цветов.

\subsection{Теорема о простом удалении вершины}

\begin{theorem}
Пусть $G$ --- неориентированный граф, $v$ --- вершина степени меньше $k$.
Если граф $G-v$ можно раскрасить в $k$ цветов, то граф $G$ также можно раскрасить в $k$ цветов.
\end{theorem}

\begin{proof}
Рассмотрим корректную $k$-раскраску графа $G-v$.
Вершина $v$ имеет меньше $k$ соседей в графе $G$, следовательно, её соседи используют не более чем $k-1$ различных цветов.
Всего доступно $k$ цветов, значит, существует хотя бы один цвет, не использованный соседями $v$.
Назначим этот цвет вершине $v$.
Тогда $v$ не конфликтует ни с одним соседом, а раскраска остальных вершин уже была корректной.
Следовательно, получена корректная $k$-раскраска всего графа $G$.
\end{proof}

\subsection{Алгоритм Чейтина--Бриггса}

\begin{algorithm}[H]
\caption{Распределение регистров через раскраску графа}
\begin{algorithmic}[1]
\Require Граф интерференции $G$, число регистров $k$
\Ensure Назначение регистров или spill
\State Построить граф интерференции по анализу живости.
\State $Stack :=$ пустой стек.
\While{$G$ не пуст}
    \If{есть вершина $v$ со степенью $<k$}
        \State удалить $v$ из $G$ и поместить в $Stack$
    \Else
        \State выбрать кандидата $v$ для spill
        \State удалить $v$ из $G$ и поместить в $Stack$
    \EndIf
\EndWhile

\While{$Stack$ не пуст}
    \State $v := pop(Stack)$
    \State выбрать для $v$ регистр, не занятый уже окрашенными соседями
    \If{такого регистра нет}
        \State пометить $v$ как spill
    \EndIf
\EndWhile

\If{есть spill}
    \State вставить загрузки и сохранения в память
    \State повторить распределение регистров
\EndIf
\end{algorithmic}
\end{algorithm}

\subsection{Минимальный пример раскраски}

Пусть есть переменные:
\[
a,b,c
\]
и граф интерференции:
\[
a-b,\qquad b-c.
\]

\[
\begin{tikzpicture}[
    node distance=1.5cm,
    v/.style={draw, circle, minimum size=0.8cm}
]
\node[v] (a) {$a$};
\node[v, right of=a] (b) {$b$};
\node[v, right of=b] (c) {$c$};

\draw (a) -- (b);
\draw (b) -- (c);
\end{tikzpicture}
\]

При двух регистрах можно назначить:
\[
a\mapsto R_1,\qquad b\mapsto R_2,\qquad c\mapsto R_1.
\]

Переменные $a$ и $c$ могут использовать один регистр, потому что они не интерферируют.

\subsection{Линейное сканирование}

Для JIT-компиляторов часто используется более простой и быстрый алгоритм --- linear scan.

\begin{algorithm}[H]
\caption{Linear Scan Register Allocation}
\begin{algorithmic}[1]
\Require Интервалы живости, отсортированные по началу
\Ensure Назначение регистров
\State $Active := \varnothing$
\For{каждый интервал $i$ в порядке возрастания начала}
    \State удалить из $Active$ интервалы, закончившиеся до начала $i$
    \If{есть свободный регистр}
        \State назначить его интервалу $i$
        \State добавить $i$ в $Active$
    \Else
        \State выбрать интервал с самым дальним концом
        \State отправить выбранный интервал в spill
    \EndIf
\EndFor
\end{algorithmic}
\end{algorithm}

Время работы обычно близко к линейному от числа интервалов.

\section{Оптимизация переходов и CFG}

\subsection{Удаление недостижимого кода}

Если блок не достижим из входа CFG, его можно удалить.

\begin{algorithm}[H]
\caption{Удаление недостижимых блоков}
\begin{algorithmic}[1]
\Require CFG $G=(V,E)$, вход $s$
\Ensure CFG без недостижимых блоков
\State Выполнить DFS или BFS из $s$.
\State Пометить все посещённые вершины как достижимые.
\State Удалить все непосещённые вершины и инцидентные им рёбра.
\end{algorithmic}
\end{algorithm}

\subsection{Упрощение переходов}

Примеры:
\[
\textbf{goto } L_1;\quad L_1:\ \textbf{goto } L_2
\quad \Rightarrow \quad
\textbf{goto } L_2.
\]

Если условие известно на этапе компиляции:
\[
\textbf{if } true \textbf{ goto } L_1 \textbf{ else } L_2
\quad \Rightarrow \quad
\textbf{goto } L_1.
\]

\subsection{Слияние блоков}

Если блок $B$ имеет единственного преемника $C$, а $C$ имеет единственного предшественника $B$, то их можно объединить.

\section{Корректность оптимизаций}

Любая оптимизация должна сохранять наблюдаемое поведение программы.

\begin{definition}
Две программы называются \textbf{семантически эквивалентными}, если для любых одинаковых входных данных они:
\begin{itemize}
    \item завершаются с одинаковым результатом;
    \item либо обе не завершаются;
    \item имеют одинаковые наблюдаемые побочные эффекты.
\end{itemize}
\end{definition}

На практике учитывают:
\begin{itemize}
    \item исключения;
    \item переполнения;
    \item неопределённое поведение языка;
    \item volatile-переменные;
    \item операции ввода-вывода;
    \item многопоточность;
    \item модель памяти.
\end{itemize}

\section{Итоговая схема оптимизации}

Типичная последовательность оптимизаций может выглядеть так:
\begin{enumerate}
    \item построение CFG;
    \item удаление недостижимого кода;
    \item вычисление доминаторов;
    \item выделение циклов;
    \item построение SSA;
    \item распространение констант;
    \item устранение общих подвыражений;
    \item удаление мёртвого кода;
    \item оптимизация циклов;
    \item межпроцедурные оптимизации;
    \item выход из SSA;
    \item распределение регистров;
    \item машинно-зависимые оптимизации.
\end{enumerate}

\end{document}
